\section{随机向量与协方差矩阵}
\label{sec:05-Probability/05-Joint-Distributions/04}

手上有十只股票。每只股票明天的收益率都是随机变量，各自的方差你知道。但拿着十个孤立的方差数字去管一个投资组合，就像你知道十个人的身高体重却不知道他们会不会同时摔倒。

两只股票。\(X_1\) 的方差 0.04，\(X_2\) 的方差 0.09。把钱平分投给两只，风险是多少？这不是两个方差平均一下的事。如果它们总是同涨同跌，你的风险比单只股票小不了多少。如果它们涨跌相反，组合波动可能远小于任何一只单独持有。

你需要的不只是方差。你需要知道每对变量之间的``共同波动''——协方差。十个变量就需要 \(10 + 45 = 55\) 个数字。怎样组织这 55 个数字，确保它们自洽，并从中读出任意方向的风险？

\subsection{把变量装进向量}

把 d 个随机变量纵向堆叠成一个列向量：

\[
X=
\begin{bmatrix}
X_1\\
\vdots\\
X_d
\end{bmatrix}.
\]

均值也自然地堆成向量：

\[
\mu = \E[X].
\]

接下来是关键。如何用一个对象同时容纳 d 个方差和 \(d(d-1)/2\) 个协方差？而且这个对象必须保证任意线性组合的方差非负——不是随便一个对称矩阵都有资格当协方差矩阵。

以下讨论假设各分量具有有限二阶矩；否则某些均值、方差或协方差可能不存在，后续矩阵运算没有意义。

\subsection{协方差矩阵的构造}

考虑 \((X-\mu)(X-\mu)^{\trans}\)。这是一个 \(d\times d\) 的随机矩阵，它的第 \((i,j)\) 元是 \((X_i-\mu_i)(X_j-\mu_j)\)。取期望：

\[
\Sigma = \Cov(X) = \E\bigl[(X-\mu)(X-\mu)^{\trans}\bigr].
\]

逐元展开看：

\[
\Sigma_{ij} = \Cov(X_i, X_j).
\]

对角线 \(\Sigma_{ii} = \Var(X_i)\)，非对角线 \(\Sigma_{ij}\)（\(i\neq j\)）是 \(X_i\) 和 \(X_j\) 的协方差。写成矩阵：

\[
\Sigma=
\begin{bmatrix}
\Var(X_1) & \Cov(X_1,X_2) & \cdots & \Cov(X_1,X_d)\\
\Cov(X_2,X_1) & \Var(X_2) & \cdots & \Cov(X_2,X_d)\\
\vdots & \vdots & \ddots & \vdots\\
\Cov(X_d,X_1) & \Cov(X_d,X_2) & \cdots & \Var(X_d)
\end{bmatrix}.
\]

协方差是对称的：\(\Cov(X_i,X_j)=\Cov(X_j,X_i)\)，所以 \(\Sigma^{\trans}=\Sigma\)。这解决了组织问题——一个对称矩阵恰好容纳所有成对信息。

但还有自洽问题。不是每个对称矩阵都能当协方差矩阵。考虑任意非随机向量 \(a\in\R^d\)，看线性组合 \(a^{\trans}X\) 这个标量随机变量的方差：

\[
\begin{aligned}
\Var(a^{\trans}X)
&= \E\bigl[(a^{\trans}(X-\mu))^2\bigr] \\
&= \E\bigl[a^{\trans}(X-\mu)(X-\mu)^{\trans}a\bigr] \\
&= a^{\trans}\,\E\bigl[(X-\mu)(X-\mu)^{\trans}\bigr]\,a \\
&= a^{\trans}\Sigma a.
\end{aligned}
\]

方差是非负的，所以对于任意向量 a，必须有 \(a^{\trans}\Sigma a \ge 0\)。协方差矩阵必须半正定（positive semidefinite）。

这个条件的几何含义是：不管沿哪个方向投影，分布的宽度不能是虚数。如果你随便写了一个对称矩阵但某个方向上的 \(a^{\trans}\Sigma a\) 是负数，那这个矩阵不可能来自任何真实的随机向量——不存在一个随机向量在某个方向上有负的方差。

\subsection{任意方向的波动：协方差矩阵的几何}

向量 a 定义了一个投影方向。标量 \(a^{\trans}X\) 是 X 沿这个方向的``影子''，它的方差

\[
a^{\trans}\Sigma a
\]

告诉你分布在这个方向上有多宽。取不同的 a，就像拿着手电筒从不同角度照一个椭球，墙上影子的宽度不同。

\(\Sigma\) 是实对称矩阵，谱定理保证它可作正交特征分解：

\[
\Sigma = Q\Lambda Q^{\trans},
\]

其中 Q 的列是互相正交的单位特征向量，\(\Lambda\) 是对角矩阵，对角线上是特征值 \(\lambda_1,\ldots,\lambda_d\)。特征向量给出分布的主轴方向，特征值给出这些方向上的方差。

如果某个特征值为零，意味着存在一个非零方向 a（对应的特征向量），使得 \(a^{\trans}X\) 的方差为零——这个线性组合几乎必然是常数。此时分布被压缩在低于 d 维的仿射子空间中：数据看似在 d 维空间里，实际只住在某个低维平面上。

这套几何不止是数学形式——PCA 的主成分、二维 Gaussian 的等密度椭圆、Kalman 滤波的误差椭球，都来自协方差矩阵的特征分解。

\subsection{线性变换怎样改变均值和协方差}

设 Y 是 X 的仿射变换：

\[
Y = AX + b.
\]

均值和协方差按简单的规则传播：

\[
\E[Y] = A\mu + b,
\]

\[
\Cov(Y) = A\Sigma A^{\trans}.
\]

第二条的推导只需要一步：\(Y-\E[Y] = A(X-\mu)\)，然后

\[
\Cov(Y) = \E\bigl[A(X-\mu)(X-\mu)^{\trans}A^{\trans}\bigr] = A\,\E[(X-\mu)(X-\mu)^{\trans}]\,A^{\trans} = A\Sigma A^{\trans}.
\]

注意是 \(A\Sigma A^{\trans}\)，不是 \(A\Sigma\) 也不是 \(\Sigma A\)。乘在哪一边由维度和对称性决定——\((X-\mu)\) 是列向量，外积是 \(d\times d\)，A 必须左乘这个外积、右乘 A 的转置才能维持维度和对称性一致。

回到投资组合的例子。两只股票收益 \(X_1,X_2\) 组成 X，投资权重向量 \(w\)，组合收益是 \(R = w^{\trans}X\)。组合方差：

\[
\Var(R) = w^{\trans}\Sigma w.
\]

展开来看：

\[
\Var(R) = w_1^2\Var(X_1) + w_2^2\Var(X_2) + 2w_1w_2\Cov(X_1,X_2).
\]

如果两只股票完全正相关（\(\Cov = \sigma_1\sigma_2\)），组合方差就是 \((w_1\sigma_1 + w_2\sigma_2)^2\)——没有任何分散效果。如果完全负相关，\(\Cov = -\sigma_1\sigma_2\)，方差变成 \((w_1\sigma_1 - w_2\sigma_2)^2\)——选对权重甚至可以精确消掉风险。分散化能否降低风险，取决于协方差矩阵的 off-diagonal 结构，而不只是单只股票的波动大小。

\subsection{白化与标准化的区别}

逐分量标准化：

\[
Z_i = \frac{X_i - \mu_i}{\sigma_i}.
\]

做完后每个分量的方差变成 1，对角线全是 1。但非对角线仍原样保留相关结构。标准化只是把每个轴向独立拉伸，相关性纹丝未动。

白化（whitening）走得更远：寻找一个线性变换 W，使得变换后的向量协方差矩阵是单位矩阵：

\[
\Cov\bigl(W(X-\mu)\bigr) = I.
\]

如果 \(\Sigma\) 正定，可以取 \(W = \Sigma^{-1/2}\)。白化消除了一切二阶线性相关——变换后任意两个不同分量的协方差都是零，每个分量的方差都是 1。

但白化不制造独立。它只处理了二阶结构。如果原向量不服从联合 Gaussian，白化后的分量可能仍存在非线性依赖（比如 U 形或尾部依赖）。只有在联合 Gaussian 的前提下，协方差为零才等价于独立——这一点和\hyperref[sec:05-Probability/05-Joint-Distributions/02]{``随机变量的独立与依赖''}一节的结论一致：独立的门槛远高于协方差层面的条件。

\subsection{协方差矩阵仍只是二阶摘要}

想一想：两个完全不同的多变量分布可以有完全相同的均值向量和协方差矩阵。一个是椭球状的单峰 Gaussian，另一个可能是几团分离的簇，或者一个甜甜圈形状的分布。它们的 \(\mu\) 和 \(\Sigma\) 一模一样，但条件依赖、尾部行为、极端事件的出现模式截然不同。

协方差矩阵擅长描述线性尺度、方向和二阶波动范围。它不能替代联合分布。在依赖关系高度非线性的场景里，盯着协方差矩阵看就像用椭圆去近似一朵云——轮廓是对了，内部结构全丢。

下一节研究一般的非线性坐标变换。线性变换时，矩和协方差的更新规则简单优美；但要得到完整密度在非线性变换下的形状，需要引入 Jacobian 来补偿局部的体积拉伸或压缩。
